\chapter{Introducción}\label{cap:introduccion}

% Norma TFGE: la página 1 de la memoria se corresponde con el inicio de
% este capítulo. La introducción debe contener síntesis del estado de la
% cuestión, objetivos del trabajo y estructura de la memoria.

\section{Contexto y motivación}

La proliferación de contenidos digitales en redes sociales, mensajería
instantánea y portales informativos ha desplazado el problema de la
desinformación desde un fenómeno puntual hasta un riesgo sistémico con
consecuencias documentadas en salud pública, procesos electorales y
estabilidad financiera \citep{lazer2018science,vosoughi2018spread}. La
detección automática de \emph{fake news}, entendida en sentido amplio como
contenido fáctualmente falso difundido con apariencia de noticia, ha
recibido en consecuencia una atención sostenida por parte de la comunidad
de aprendizaje automático y procesamiento del lenguaje natural (PLN).

Históricamente, el problema se ha planteado como una clasificación binaria
\textsc{real}/\textsc{falso} sobre el cuerpo de la noticia, utilizando
representaciones léxicas (TF-IDF, $n$-gramas) y modelos lineales o
ensembles \citep{conroy2015automatic,shu2017fake}. Esta formulación, sin
embargo, presenta limitaciones de fondo:

\begin{enumerate}[leftmargin=*,itemsep=0.2em]
  \item El veredicto se aplica al texto completo, lo que ignora que un
        artículo puede combinar afirmaciones verdaderas, falsas y otras
        sin evidencia suficiente.
  \item La etiqueta es opaca para el usuario: un dictamen «falso» sin
        justificación verificable carece de utilidad operativa y de
        garantías de auditoría.
  \item El modelo aprende correlaciones de estilo o de fuente que no se
        generalizan fuera del dominio de entrenamiento.
\end{enumerate}

A partir de 2018, el dataset \textsc{FEVER} (\emph{Fact Extraction and
VERification}) \citep{thorne2018fever} reformula el problema como
\emph{verificación de afirmaciones contra evidencias}: dado un enunciado
$c$ (\emph{claim}) y un conjunto de evidencias $\mathcal{E} = \{e_1,
\dots, e_n\}$ extraídas de una base de conocimiento, el sistema debe
predecir uno de tres veredictos: \textsc{supported}, \textsc{refuted} o
\textsc{not-enough-info}. Este planteamiento alinea la tarea con la
inferencia de lenguaje natural (NLI) y permite separar dos cuestiones
estadísticas distintas: (i) la calidad del clasificador NLI por par
$(c, e)$ y (ii) la regla de agregación que combina los veredictos
parciales en un dictamen global.

\section{Objetivos}

El presente trabajo aborda la componente \textbf{estadística} y de
\textbf{modelización predictiva} de un sistema de verificación de
afirmaciones, complementario a la componente de ingeniería de software
descrita en una memoria independiente del Grado en Ingeniería Informática
del mismo autor. Los objetivos específicos son:

\begin{enumerate}[leftmargin=*,itemsep=0.2em]
  \item Formalizar la verificación de afirmaciones como un problema de
        clasificación condicional $\Prb(Y = y \mid c, e)$ con
        $y \in \{\textsc{sup}, \textsc{ref}, \textsc{nei}\}$, y derivar
        de él una regla de decisión a nivel de \emph{claim} con cuatro
        posibles veredictos (añadiendo el caso \textsc{conflicting}
        cuando las evidencias se contradicen).
  \item Implementar y evaluar un \textbf{baseline} basado en
        representaciones TF-IDF y regresión logística multinomial, con
        regularización $\ell_2$, validado mediante \emph{stratified
        $k$-fold} y diagnóstico de calibración.
  \item Realizar el \emph{fine-tuning} de un modelo neuronal
        \textsc{DeBERTa-v3} sobre FEVER y comparar su rendimiento con el
        baseline en términos de exactitud, F1 macro, log-loss y
        \emph{Expected Calibration Error} (ECE).
  \item Aplicar técnicas de \textbf{calibración \emph{post-hoc}} (Platt e
        isotónica) sobre las probabilidades del clasificador y cuantificar
        su efecto mediante diagramas de fiabilidad.
  \item Realizar \textbf{contrastes estadísticos} formales para comparar
        modelos: test de McNemar para diferencias en error de
        clasificación pareadas, y \emph{bootstrap} no paramétrico para
        construir intervalos de confianza sobre la diferencia de F1
        macro.
  \item Proponer y analizar una \textbf{regla de agregación} probabilística
        que combine los veredictos por evidencia en un veredicto por
        \emph{claim}, discutiendo sus propiedades (consistencia,
        comportamiento bajo evidencias contradictorias, manejo de
        ausencia de información).
\end{enumerate}

\section{Estructura de la memoria}

El capítulo~\ref{cap:estado-arte} revisa el estado del arte en detección
de noticias falsas y verificación de afirmaciones. El
capítulo~\ref{cap:marco-teorico} desarrolla el marco teórico-estadístico:
clasificación multinomial, regresión logística, modelos de atención y la
formulación NLI. El capítulo~\ref{cap:diseno} describe el diseño
experimental sobre FEVER (particiones, métricas y protocolo de
evaluación). Los capítulos~\ref{cap:baseline} y~\ref{cap:deberta}
desarrollan respectivamente el baseline lineal y el modelo neuronal. El
capítulo~\ref{cap:calibracion} aborda la calibración \emph{post-hoc}. El
capítulo~\ref{cap:tests} formaliza los contrastes estadísticos
empleados. El capítulo~\ref{cap:agregacion} estudia la regla de
agregación de claims. Finalmente, el capítulo~\ref{cap:resultados}
presenta y discute los resultados, y el capítulo~\ref{cap:conclusiones}
recoge las conclusiones y posibles líneas de trabajo futuro.

Las salidas de ordenador completas, configuraciones de entrenamiento y
fragmentos de código fuente se entregan en un fichero separado de
anexos, conforme a la normativa del TFG de Estadística.
